\section{Efficient pairs cannot keep repeating}\label{sec:pairs}

We begin with the fourteen-symbol case, because the improvement can be
seen before introducing its general parameters. Here $n=13$. A chain
has length at most $36$, and in a path with two or more chains its length
is at most $35$. The earlier local collision restrictions say more: a
chain longer than four forces its successor to have length at most four,
or at most three if that successor is not the last chain in the path.

Pair each nonterminal chain longer than four with its successor. The
pairs are disjoint, because the successor is too short to begin another
pair. All remaining chains, except possibly the last, have length at
most four. An interior pair has total length at most $38$, attained
only by $(35,3)$. Thus the earlier count allows an average of $19$
regular clusters per chain. The length count by itself would permit
$(35,3),(35,3),\ldots$ indefinitely.

Now follow what such a pair does to the symbols. Its normalization has
nonzero cycles
\[
 (1\ 11\ 7\ 3\ 9),\qquad(2\ 10\ 5\ 12\ 6),\qquad(4\ 13\ 8).
\]
After fifteen pairs the relabeling is the identity. If another chain
then starts, it starts in the same rotation class as the first one.
Consequently fifteen consecutive interior pairs cannot all be $(35,3)$.
At least one must have smaller total length, or the pairing must be
interrupted by a singleton. Either event loses at least one regular
cluster from the old allowance.

There are at most two chains per pair or singleton. Thus one lost unit
per fifteen such pieces predicts a reduction of $1/30$ in the mean
chain length. Equivalently, the parameter $t=38$ should decrease by
$1/15$. This is the source of the new value $569/15$. The argument
below proves this prediction while allowing for the final chain, whose
successor need not exist.

\subsection{Counting the interruptions}

Call a singleton or a chosen pair a \emph{tile}. If the old mean chain
length is $\mu$, a tile containing $k$ chains is allowed total length
$\mu k$. Its \emph{loss} is the difference between this allowance and
its actual total length. This loss counts regular clusters; the run
deficits introduced earlier count missing blocks.

The count is simple once the efficient tiles have been identified.
Mark every tile with positive loss. These marks divide the zero-loss
tiles into consecutive strings. A bound on the lengths of those strings
then forces enough marks. The only endpoint issue is that the last tile
may attain a larger length, and a return after it supplies no contradiction
unless a further chain actually starts. We therefore set that tile aside.

\begin{lemma}\label{lem:tiles}
Suppose a path of $q\ge2$ chains is partitioned into tiles of one or
two chains, and a tile with $k$ chains and total length $L$ has loss
$\mu k-L$. Suppose that before the last tile every loss is nonnegative,
every positive loss is at least $\varepsilon>0$, and every consecutive
string of zero-loss tiles has length at most $R-1$, where $R\ge1$.
If the last tile has loss at least $-c$, where $c\ge0$, then
\[
 \sum_{i=1}^q\ell_i\le
 \left(\mu-\frac{\varepsilon}{2R}\right)q+c+\varepsilon.
\]
\end{lemma}

\begin{proof}
Let $B_+$ be the number of marked tiles and $P$ their total loss,
excluding the last tile. Then $B_+\le P/\varepsilon$. There are at
most $B_++1$ intervening strings, each with at most $R-1$ zero-loss
tiles. Restoring the last tile, the total number of tiles is at most
$R(B_++1)$. Since each tile has at most two chains,
\[
 q\le2R(P/\varepsilon+1),\qquad
 P\ge\frac{\varepsilon q}{2R}-\varepsilon.
\]
The total chain length is at most $\mu q-P+c$, giving the result.
\end{proof}

\subsection{The case \texorpdfstring{$n=3m+1$}{n=3m+1}}

The fourteen-symbol example is typical of this residue class. Assume
$m\ge3$, put $p_0=m(2m+1)$ and $b=p_0-1$, and pair every nonterminal
chain longer than $m$ with its successor. Lemma~\ref{lem:inherited-chains} bounds that successor by $m$, and
by $m-1$ when it is nonterminal. Hence an interior pair has total
length at most $t_0=p_0+m-2$, with equality only for $(b,m-1)$.

Take $\mu=t_0/2$. An interior singleton has length at most $m$, so
its loss is at least $\mu-m\ge1$. All positive interior losses are
integers, since $t_0=2m^2+2m-2$ is even. Thus a tile either is the
efficient pair $(b,m-1)$ or loses at least one unit.

The normalization of that pair is
\begin{equation}\label{eq:efficient-map}
 M_1=h_{m-1}h_b=\rho g^{m-1}\rho g^{-1}.
\end{equation}
Let $R_1$ be its order. As in the example, $R_1$ consecutive efficient
interior tiles would return to the first start class. The next start
is forced because a deficit-two run has $n-2\ge8$ blocks, above the
three-block threshold after $B$.

Merely knowing that $M_1$ has finite order would give a weak quantitative
estimate. Here its cycles can be calculated explicitly. Their lengths
are listed in Table~\ref{tab:cycles}, and the calculation in
Appendix~\ref{app:cycles} proves $R_1\le n^2/4$. The example $n=13$
has cycle lengths $5,5,3$, so $R_1=15$.

It remains to account for the last tile. A final singleton can have
length $b$, giving loss $-(p_0-m)/2$. A final pair can have total
length $b+m=t_0+1$, giving loss $-1$. Thus the single allowance
$c=(p_0-m)/2$ covers both possibilities. Apply Lemma~\ref{lem:tiles}
with $\varepsilon=1$. Since $p_0-\mu=c+1$, it yields
\[
 \sum_i\ell_i\le
 \left(\mu-\frac1{2R_1}\right)q+c+1
 \le p_0+\frac12\left(t_0-\frac1{R_1}\right)(q-1).
\]
The case of a single chain already has the stronger bound $p_0$.
We have therefore proved
\begin{equation}\label{eq:CP1}
 \CP\left(p_0,t_0-\frac1{R_1}\right).
\end{equation}
At fourteen symbols, the final-tile allowance is $16$. It remains an
endpoint allowance; it does not cancel the losses that accumulate along
a long path.

\subsection{When the interiors collide first}

For $n=3m+2$, following the first class of each tile misses an earlier
collision. It is useful to recall the full run in Figure~\ref{fig:run}.
Delete the symbol zero from each of its four rotation classes. All four
give the same cyclic order $1234$; they differ only in the gap into which
zero was inserted. In general, a full run with distinguished symbol zero
visits all $n$ such insertions into one cyclic order of the other symbols.
This family of classes is called a \emph{two-loop}.

Consequently two full runs with the same nonzero cyclic order repeat
every class, even if the starts of the larger pieces containing them
are different. This is what happens in the efficient pattern for the
present residue. Put $p_0=(m+1)(2m+1)$, $v=2m(m+1)$, $b=p_0-1$,
and $c=v-2$.

\begin{lemma}\label{lem:two-pairs}
For $m\ge2$, the pattern $T_bT_cT_bT_c$ cannot occur in a genuine
light segment.
\end{lemma}

The repeated runs lie inside the two shorter chains. To locate them,
look at the $C$ cells of each $T_c=C^{c-1}B$. In the first copy choose
the cell at distance $m$ from the end of this $C$ portion; in the second
choose the cell at distance $m$ from its beginning, as in
Figure~\ref{fig:collision}. The last full run of the first chosen cell
and the first full run of the second chosen cell have the same nonzero
cyclic order. Appendix~\ref{app:full-collision} checks this equality by
following the relabelings. Both runs lie strictly before their chains'
final $B$ cells, so the collision occurs within the displayed pattern.

\begin{figure}[htbp]
\centering
\renewcommand{\arraystretch}{1.6}
\begin{tabular}{r l l}
first pair & $T_b\quad C^{c-2-m}\,\boxed{C}\,C^mB$
 & last full run in the box\\
second pair & $T_b\quad C^m\,\boxed{C}\,C^{c-2-m}B$
 & first full run in the box
\end{tabular}
\caption{The two full runs that coincide. The boxes are at equal depths
from opposite ends of the $C$ portions of the shorter chains. Their
cyclic orders are compared in the coordinates of the whole path.}
\label{fig:collision}
\end{figure}

We can now use the same loss count. Lemma~\ref{lem:inherited-chains} bounds the successor of a chain
longer than $v$ by $v-1$. If the first length is at least $v+2$ and
the successor is nonterminal, the stronger bound is $v-2$.
Pair every nonterminal chain longer than $v$ with its successor, and
put $t_0=b+v-2=4m^2+5m-2$ and $\mu=t_0/2$.

For $m>2$, the only interior pair attaining $t_0$ is $(b,v-2)$.
A pair beginning with $v+1$ loses at least $m-2$, and every other
positive pair loss is an integer. A singleton loses at least
$(m-2)/2$. Thus every positive interior loss is at least
$\varepsilon=\min\{1,(m-2)/2\}$. Lemma~\ref{lem:two-pairs}
excludes two successive efficient tiles, so take $R=2$.

A final singleton has loss at least $-(m+2)/2$, and a final pair has
loss at least $-1$. Lemma~\ref{lem:tiles}, together with
$p_0-\mu=(m+2)/2+1$ and $\varepsilon\le1$, gives
\begin{equation}\label{eq:CP2}
 \CP(p_0,t_0-\varepsilon/2),\qquad m>2.
\end{equation}
The decrement is $1/4$ at $m=3$ and $1/2$ thereafter.

At $m=2$ there are three efficient tiles: $(12)$, $(13,11)$ and
$(14,10)$. Their normalizations lie in a group of order $12$, so
twelve consecutive efficient interior tiles would have two equal
prefix normalizations and hence a repeated start class. Positive
losses are integers and the final allowance is two. The same lemma
gives $\CP(15,287/12)$. The small group calculation is also included
in Appendix~\ref{app:full-collision}.

All these collisions occur within the path under consideration. Thus
the counts apply in every zero-ended factor, as required by $\CP$;
no completion of such a factor to a superpermutation is assumed.
